%% Part I -- Intro (comic pages 6 to 10).
%% The section that establishes why the book exists: the fear inherited from
%% fiction, the two authors and what each is after, and the turn from
%% replication to augmentation that the rest of the book runs on.
%%
%% ONE PANEL CAPTION, ONE SLIDE. The comic sets every topic page as six
%% panels in a 2x3 grid with one caption under each. Each caption takes a
%% slide of its own so the room reads one claim at a time; the six-dot row
%% under it marks which panel is on screen, and six consecutive slides
%% sharing a frame title are one comic page, read in the comic's own order.
%% Captions are verbatim, checked against a render of each page.
%% \setpagethesis opens each run with that page's one line, mine, not theirs.
%% Fragment only: no preamble, \input by ai_for_all.tex.

\sectionsubtitle{pages 6 to 10}
\section{Intro}

% Page 6 is the one topic page the comic does NOT set as six panels: a single
% line across an otherwise empty spread. So it keeps the single-passage
% \comicquote treatment while every page after it runs as six slides.
\begin{frame}{Knowing what future we would like to have}
  \panel{page 6, the opening spread}
  \begin{comicquote}{6}
    Knowing the future is impossible. But what we can do: Knowing what
    future we would like to have. And then work on it.
  \end{comicquote}
\end{frame}

% The page that sets the posture. Golem and Frankenstein establish that the
% fear is old, the rubber boots make it domestic, and Hawking turns it into a
% question with a deadline. Panels 5 and 6 are one Hawking quotation split
% across two panels.
\setpagethesis{The fear is older than the technology. Hawking only gave it a deadline.}
\begin{frame}{Our creatures have always turned against us}
  \panel{page 7, Golem, Frankenstein and the rubber boots}
  \panelslide{7}{1}{Our creatures have always turned against us. This applies to the Golem of Jewish mythology as well as to Frankenstein.}
\end{frame}
\begin{frame}{Our creatures have always turned against us}
  \panel{page 7, Golem, Frankenstein and the rubber boots}
  \panelslide{7}{2}{Neither magic nor natural sciences have ever helped to make artificial humans submissive and controllable.}
\end{frame}
\begin{frame}{Our creatures have always turned against us}
  \panel{page 7, Golem, Frankenstein and the rubber boots}
  \panelslide{7}{3}{Parents whose kids prefer to wear rubber boots instead of weather-proof sandals even when it's hot outside know that.}
\end{frame}
\begin{frame}{Our creatures have always turned against us}
  \panel{page 7, Golem, Frankenstein and the rubber boots}
  \panelslide{7}{4}{An AI Take-Over is a hypothetical scenario in which AI takes control of the planet -- taking it away from us.}
\end{frame}
\begin{frame}{Our creatures have always turned against us}
  \panel{page 7, Golem, Frankenstein and the rubber boots}
  \panelslide{7}{5}{In 2011, Stephen Hawking said that `Success in creating AI would be the biggest event in history.'}
\end{frame}
\begin{frame}{Our creatures have always turned against us}
  \panel{page 7, Golem, Frankenstein and the rubber boots}
  \panelslide{7}{6}{`Unfortunately, it might also be the last, unless we learn how to avoid the risks.' How can we make AI systems both safe and beneficial?}
\end{frame}

% Julia, the economist. Panel 3 is the bias claim the whole Risks section
% later rests on, and it is made as a compliment to data rather than an
% accusation.
\setpagethesis{Data will not lie to you strategically. It will still lie to you.}
\begin{frame}{Who is Julia}
  \panel{page 8, Julia, the digit wall and the 404}
  \panelslide{8}{1}{Good question. That's something I ask myself, now and then.}
\end{frame}
\begin{frame}{Who is Julia}
  \panel{page 8, Julia, the digit wall and the 404}
  \panelslide{8}{2}{I love making sense of our world. I love to discover patterns in data, with code.}
\end{frame}
\begin{frame}{Who is Julia}
  \panel{page 8, Julia, the digit wall and the 404}
  \panelslide{8}{3}{I love data's friendliness, their lack of strategic answering. If they are biased it's not their fault but ours.}
\end{frame}
\begin{frame}{Who is Julia}
  \panel{page 8, Julia, the digit wall and the 404}
  \panelslide{8}{4}{Our brains are prone to some errors. Overconfidence. And discrimination against people who seem to be different from us.}
\end{frame}
\begin{frame}{Who is Julia}
  \panel{page 8, Julia, the digit wall and the 404}
  \panelslide{8}{5}{Artificial Intelligence (AI) might have the potential to help us find answers to problems we are not able to solve.}
\end{frame}
\begin{frame}{Who is Julia}
  \panel{page 8, Julia, the digit wall and the 404}
  \panelslide{8}{6}{For that reason, we need to demystify AI. Therefore, I'm writing this comic. Let your voice be heard.}
\end{frame}

% Lena, the illustrator, and the answer to why this is a comic and not a
% paper: the subject is a mirror, so the medium is one too.
\setpagethesis{Ask what the way we draw AI says about what we expect from it.}
\begin{frame}{Who is Lena}
  \panel{page 9, Lena, the iceberg of visual codes}
  \panelslide{9}{1}{Not one thing in particular, I guess.}
\end{frame}
\begin{frame}{Who is Lena}
  \panel{page 9, Lena, the iceberg of visual codes}
  \panelslide{9}{2}{I love visual language. I love to encrypt thoughts into pictures by exploiting society's visual memory.}
\end{frame}
\begin{frame}{Who is Lena}
  \panel{page 9, Lena, the iceberg of visual codes}
  \panelslide{9}{3}{Visuality is one of the most vivid encryption systems I know. We are constantly increasing and renewing it collectively.}
\end{frame}
\begin{frame}{Who is Lena}
  \panel{page 9, Lena, the iceberg of visual codes}
  \panelslide{9}{4}{It's fun to mess around with visual codes by recomposing associations. Pictures start to shift meanings and interfere in society's processes.}
\end{frame}
\begin{frame}{Who is Lena}
  \panel{page 9, Lena, the iceberg of visual codes}
  \panelslide{9}{5}{AI is a field that gives us an intimate view on society's desires and fears.}
\end{frame}
\begin{frame}{Who is Lena}
  \panel{page 9, Lena, the iceberg of visual codes}
  \panelslide{9}{6}{Looking from different angles might help us to overcome commonly known perspectives. Let your view be seen!}
\end{frame}

% Maschinen-Maria 1927, R2D2 1977, T-800 1984, Samantha 2013, Tars 2014.
% Five panels of fiction and one panel that turns the parade into an argument.
\setpagethesis{Every AI in the parade is a character. Not one of them is a system.}
\begin{frame}{Lessons from AI in Fiction}
  \panel{page 10, the five fictional robots, 1927 to 2014}
  \panelslide{10}{1}{}
\end{frame}
\begin{frame}{Lessons from AI in Fiction}
  \panel{page 10, the five fictional robots, 1927 to 2014}
  \panelslide{10}{2}{}
\end{frame}
\begin{frame}{Lessons from AI in Fiction}
  \panel{page 10, the five fictional robots, 1927 to 2014}
  \panelslide{10}{3}{}
\end{frame}
\begin{frame}{Lessons from AI in Fiction}
  \panel{page 10, the five fictional robots, 1927 to 2014}
  \panelslide{10}{4}{}
\end{frame}
\begin{frame}{Lessons from AI in Fiction}
  \panel{page 10, the five fictional robots, 1927 to 2014}
  \panelslide{10}{5}{}
\end{frame}
\begin{frame}{Lessons from AI in Fiction}
  \panel{page 10, the five fictional robots, 1927 to 2014}
  \panelslide{10}{6}{Seems, like we should not seek to replicate but to augment ourselves.}
\end{frame}
